\begin{theorem}[Wild quadratic degree-two coordinate transport]
\label{mod:cases-wildquadratic-coordinatetransport}
Provide the common degree-two coordinate and phase-transport interface used
by every positive-polar Wild quadratic branch.  The public proposition
\texttt{WildQuadraticPositivePolarCoordinateTransport} records, uniformly in
the refinement parameter \(q\), the correction-class provenance, the exact
upper coordinate certificate, and equality of the actual extension
Lamprecht factor with the coefficient-table phase.  Its companion
continuation retains the lower product coordinate and the actual
extension, norm, base, and twist phase equalities without choosing \(q\).

The node contains only shared norm-polynomial, normalized-ratio, critical-unit,
coordinate-point, and actual phase-source transport.  Low, boundary, and
strict-high source mathematics is supplied by the three downstream producer
nodes.
\LeanFile{LanglandsFirstMainLemma/Cases/WildQuadratic/CoordinateTransport.lean}
\lean{LanglandsFirstMainLemma.WildQuadraticPositivePolarCoordinateTransport}
\leanok
\SourceText{Lemma lem:critical-polar-coordinate; Lemma lem:residual-parameter-transport; Propositions prop:quadratic-critical-functions, prop:quadratic-gamma-table, and prop:quadratic-complete-assembly.}
\uses{mod:parameters-phasereduction,mod:parameters-highquadratic,mod:cases-wildquadratic-coefficienttable,mod:cases-wildquadratic-errorformula}
\FormalizationContract
The interface is polymorphic in \(q\).  Every stationary numerator remains a
quotient class until a source-specific theorem supplies its representative.
Elementary stationary values and residual Hasse phases are transported
together, so the displayed equalities concern actual Lamprecht factors, not
merely coefficient pairs.  No low or high source representative is selected
in this shared node.
\end{theorem}
